\scanpage{111}
\mspage{54}

\deleted{sur lequel $\Sigma$ opère} \inserted{et l'opération de $\Sigma$} \deleted{\unsure{dessus}} sur $\pi$ est
compatible avec cette structure à lacets,
et le caractère cyclotomique $\chi:\Gamma\to\widehat{\mathbb Z}^{*}$
(\unsure{plutôt}, $\chi\mid\Sigma$). Ainsi l'opération de
$\Sigma$ sur $\pi$ implique son action sur $I$, d'où
$S(\overline K)\simeq\unsure{I_0|\overline S_0}\wedge_{\Sigma}\Gamma$ en tant que $\Sigma$-ensemble
--- on récupère donc le $K$-schéma étale $S$.

Mais mieux, on récupère tout le
diagramme
\diagram{\begin{tikzcd}[ampersand replacement=\&]
\Pi_{\overline D_0^{*}} \arrow[r] \& {[\Pi_{\overline U_0}\unsure{\to|-}]}\,\Pi_{\overline U} \arrow[r] \& \Pi_e
\end{tikzcd}}
et l'opération de $\Sigma$ dessus, d'où le
diagramme des topos \unsure{des revêt.}\inserted{\deleted{``\unsure{étales}''}\illegible{1}} --- ou si
on préfère, le diagramme
\diagram{\begin{tikzcd}[ampersand replacement=\&]
\Pi_{D^{*}} \arrow[r] \& \Pi_U \arrow[r] \& \Pi_e
\end{tikzcd}}
(avec les notations du début de ces notes,
qui deviendraient ici $\Pi_{D^{*},\Gamma}\to\Pi_{U,\Gamma}\to\Pi_{e,\Gamma}$),
\unsure{plus} \unsure{bien|rien} sur $K$\dots

Les hom $(X',S')\xrightarrow{f}(X,S)$
[$S'\supset f^{-1}(S)$, $f$ dominant] \deleted{\unsure{correspondent}} \inserted{se décrivent}
\deleted{\unsure{simplement aux}} \inserted{simplement} (via le choix d'un $\overline X'_0$
au dessus d'un $\overline X_0$) \deleted{aux} \inserted{par des} hom extérieurs

\margin{\uline{\unsure{(159)}}}
\margin{110}

\scanpage{112}

$\pi'\deleted{\illegible{1}}\ (=\pi_1(\overline X'_0))\longrightarrow\pi\ (=\pi_1(\overline X_0))$, compatibles
avec les actions de $\Sigma'\ (\subset\Sigma)$ et de $\Sigma$,
et avec les structures de lacets \unsure{anabéliens}
(ce qui s'exprime \deleted{par} \inserted{à l'aide d'} une application
$\tau:I'\deleted{\unsure{\subset}\overline S'_0\illegible{1}}\longrightarrow I=\overline S_0$ compatible avec
$\Sigma'$, et une application d'entiers naturels
$(d_i)_{i\in I'}$).
\margin{\uline{NB} $I'=\overline S'_0\cap f^{-1}(\overline S_0)$}

Il faut cependant compléter, \deleted{\illegible{1}} pour
$I=\varnothing$ i.e. $S=\varnothing$, la définition de la
structure à lacets, par la donnée d'un
isom
\[
\kappa:T^{\otimes-1}\simeq H^2(\pi,\widehat{\mathbb Z})
\]
[NB, en caractéristique $p>0$, on doit se borner
aux composantes $l$-adiques avec $l\ne p$]
i.e.
\[
\widehat{\mathbb Z}\simeq H^2(\pi,T)
\]
compatible avec l'action de $\Sigma$
--- et il faut exiger, dans l'interprétation
``galoisienne'' de $f:(X',S')\unsure{\to}(X,S)$, quand $S=S'=\varnothing$
que l'hom $\pi\to\pi'$ induit un
diagramme commutatif

\margin{111}

\scanpage{113}
\mspage{55}

\diagram{\begin{tikzcd}[ampersand replacement=\&]
\widehat{\mathbb Z} \arrow[r, no head, "\unsure{\kappa^{*}|\kappa}", "\sim"'] \arrow[d, "d"'] \& H^2(\unsure{\pi},T) \arrow[d, "f^{*}"] \\
\widehat{\mathbb Z} \arrow[r, "\kappa'", "\sim"'] \& H^2(\pi',T)
\end{tikzcd}},
où $d\inserted{\in\mathbb N}$ est le \uline{degré} (défini \deleted{\illegible{2}} \inserted{de façon unique}
par cette condition,
\deleted{galoisienne} \inserted{ou aussi} l'ordre de $\operatorname{Coker}f^{*}$). \deleted{\illegible{1}}

\deleted{à priori on aurait donc} \inserted{Quand $\unsure{S^{*}|S}=\varnothing$, $\unsure{S'|S'_0}\ne\varnothing$, il devrait y avoir}
encore y avoir une compatibilité
pour les données \unsure{harmoniciennes} \unsure{(de}
nature différente sur $\pi'$, où il y a bel et
bien ``des lacets'', et sur $\pi$, où elle est
purement cohomologique. \unsure{\S} La question
\illegible{1} sans doute : celle de décrire
une structure \unsure{harmonicienne} ``cohomologique''
sur le groupe $\widetilde\pi'$, déduit d'un $\pi'$ à
lacets (avec $I'\ne\varnothing$) en divisant par
lesdits lacets. La question est la
même, semble-t-il, dans le cadre topologique,
ou le cadre arithmétique, il me faudra
revenir dessus. Il faudrait que pour \unsure{tt}

\margin{\unsure{paradigme du passage de} $(X,S)$ \unsure{à} $X$ : ``bouchage des trous''\dots}
\margin{112}

\scanpage{114}

hom \inserted{$\pi'\to\pi$} de groupes à $T$-lacets, l'hom
$\unsure{\widehat\pi'|\widetilde\pi'}\to\widetilde\pi$ correspondant respecte aussi (pour
un degré convenable) la structure à lacets.

Quand on s'intéresse aux \unsure{syst.} anabéliens
$(\overline X,\overline S)$ définis directement \uline{sur $\overline K$}, avec $\overline X$
connexe disons, ceci s'exprime
par un groupe extérieur à $T$-lacets,
muni d'une action non de $\Gamma$ mais
du \uline{\unsure{noyau|germe}} de groupe profini défini par $\Gamma$.
Si on considère les \deleted{\illegible{1}} ``$\Gamma$-\unsure{autom}'' d'un
tel objet (formés d'un $\gamma\in\Gamma$ et
d'un autom \inserted{ext.} $f$ de $\pi$ respectant la
structure à lacets, $f$ et $\gamma$ étant compatibles
dans un sens évident\dots), on trouve
un groupe \inserted{\deleted{profini} (pro-fini ???)} $G$ et un hom \inserted{(``discret'' ???)} $G\to\Gamma$
à image un ss-groupe ouvert de $\Gamma$, et
à noyau \inserted{$G_0$} le groupe des autom \unsure{gal.|ext.} à
lacets de $\pi$ qui commutent à l'action du
``\unsure{noyau|germe}'' \deleted{\unsure{typique}} (lequel $G_0$ est conjecturalement par la
``\unsure{conj. bordélique}'', isom à \inserted{au gr fini} $\operatorname{Aut}_{\overline K}(\overline X,\overline S)$\dots).

\margin{113}

\scanpage{115}
\mspage{56}

En termes de cette suite exacte
\[
1\unsure{\longrightarrow|-}G_0\longrightarrow G\unsure{\longrightarrow|-}\Gamma
\]
la ``\unsure{restriction|restitution}'' de $(\overline X,\overline S)$ au corps $K$ s'exprime
donc (on l'espère, de façon pleinement fidèle,
si la conj. \unsure{bordélique} est valable) par
un scindage $\Gamma\to G$ de $G\to\Gamma$\dots

$13^{\mathrm{bis}}$. \uline{Retour sur la notion de groupe à lacets}

Soit $\pi$ un groupe, $([L_i])_{i\in I}$ une famille
de classes de conjugaison de ss-groupes de
$\pi$. On dit que ça définit une \deleted{\unsure{des}}
``structure à lacets'' sur $\pi$ \inserted{[de type $(g,\nu)$]}, si $\exists\,g\in\mathbb N$,
$\unsure{\forall i\in I}$ un $L_i\in[L_i]$, un générateur
$l_i\in L_i$, des éléments $\unsure{x_\alpha},y_\alpha\in\pi$ $(1\le\alpha\le g)$
enfin un ordre $i_1\dots i_\nu$ sur $I$ (\deleted{de} \inserted{$\nu=$} \unsure{cardinal|card} $I$),
tels que (posant $l_\alpha=l_{i_\alpha}$ pour simplifier)
\[
[x_1,y_1][x_2,y_2]\dots[x_g,y_g]l_1\dots l_\nu=0
\]
sont une présentation du groupe $\pi$. On
n'exclut pas a priori ni le cas $g=0$, ni le cas $\nu=0$,
i.e. $I$ vide.

On déduit de ceci :

a) Si $\nu\ne0$, $\pi$ est \uline{libre} (à $2g+\nu-1$ générateurs) ---
donc libre \uline{non abélien} sauf si $g=0$, $\nu\le2$ \deleted{\illegible{1}}.

\margin{114}

\scanpage{116}

b) Si $\nu=0$, le seul cas où $\pi$ abélien est celui
$g\le1$. [Donc $\pi\simeq\Pi_{g,\nu}$ est abélien ssi
$g=0$, $\nu\le3$ ou $g=1$, $\nu\unsure{\le|=}0\deleted{\unsure{ou 1}}$]
\margin{\unsure{Bien que} $\pi=0$ \unsure{cas} \deleted{\illegible{1}} $\unsure{g=0}$, $\nu=0,1$, \illegible{3} \unsure{en ce cas} \illegible{2}.}

En tous cas, on a une suite exacte
canonique de $\mathbb Z$-modules libres de t.f.
\[
0\unsure{\longrightarrow|-}T_\pi\longrightarrow\prod_{i\in I}L_i
\xrightarrow{i}\pi_{\mathrm{ab}}\unsure{\longrightarrow|-}\widetilde\pi_{\mathrm{ab}}\longrightarrow0
\]
\margin{\uline{NB} Sauf si $\pi=0$, i.e. $g=0$, $\nu=0,1$, les $L_i$ sont $\simeq\mathbb Z$.}
où, \inserted{pour $I\ne\varnothing$, $\pi\ne0$} $T_\pi$ est $\simeq\mathbb Z$ (défini comme $\operatorname{Ker}i$), et où
les projections
\[
T_\pi\longrightarrow L_i
\]
sont des iso. Dans le cas $I\ne\varnothing$, on
appelle $T_\pi$ le \uline{$\mathbb Z$-module des orientations de}
$\pi$ (muni de la famille des $(L_i)$) ---
on définit, \deleted{plus \illegible{1} $T_\pi$ \unsure{si} $I$ \illegible{2}} \inserted{si $I=\varnothing$, $(\nu=0)$, mais $g\ne0$}
(donc $\pi\ne1$ --- le seul cas $\pi=1$ est celui où
$g=0$, $\nu=0$ ou $1$) par la condition
\[
\unsure{T_\pi}=\underbrace{H^2(\pi,\mathbb Z)}_{\substack{\mathbb Z\text{-module}\\\text{libre de rg }1}}{}^{\otimes-1}
\]
[on établira plus bas une relation entre les
deux définitions de $T_\pi$]. Ainsi $\pi_{\mathrm{ab}}$
est \inserted{libre} de rang $2g+\nu-1$ si $\nu\ne0$, $2g$ si $\nu=0$.
Donc $g$ est uniquement déterminé par $\pi$ et $\nu=\operatorname{card}I$.
\deleted{la famille des $([L_i])$.}

\margin{115}

\scanpage{117}
\mspage{57}

Notons que [si $\nu\ne0$, \inserted{et]} sauf les cas ``abéliens''
$g=0$, $\nu=1,2$, les classes de conjugaison
des $L_i$ $(i\in I)$ sont distinctes, donc la
structure à lacets de $\pi$ peut se décrire
comme la donnée d'un ens. de $\nu$ classes de
conjugaison de ss-groupes de $\pi$. De plus,
\deleted{(du fait que $L_i\to\pi_{\mathrm{ab}}$ injectif}
on voit que $\forall\,g\in\pi$ qui normalise
un $\unsure{L_i}$ le centralise (c'est évident
en tous cas pour $\nu\ge2$, car alors
$L_i\to\pi_{\mathrm{ab}}$ est injectif), ce qui implique
que les $\{L_i\}$ d'une même classe $[L_i]$
sont canoniquement isomorphes entre eux
\deleted{(\illegible{2} :}
ce qui donne un sens intrinsèque au
terme $\prod L_i$ dans la suite exacte
plus haut, et à l'iso. can. $T_\pi\to L_i$\dots
\margin{Voir à part le cas $\nu=1$ : normalisateur du ss-groupe engendré par $\prod[x_\alpha,y_\alpha]$, dans le groupe libre engendré par les générateurs $x_\alpha,y_\alpha$\dots}

Si un groupe $\Sigma$ opère sur la structure
à lacets $(\pi,([L_i]))$ il opère sur le
$\mathbb Z$-module inversible $T_\pi$, d'où un
caractère
\[
\chi:\Sigma\longrightarrow\mathbb Z^{*}=\{\pm1\}
\]
---
\margin{On suppose dorénavant qu'on est dans le cas anabélien ; \unsure{pour le moins} on exclut $g=0$, $\nu\le2$.}

\margin{116}

\scanpage{118}

inversement, si on a un caractère \inserted{$\chi$ sur $\Sigma$} donné d'avance,
on parlera de \deleted{\illegible{1}} action de $\Sigma$ sur $(\pi,([L_i]))$
\uline{compatible} avec $\chi$. Si on se donne un
$\mathbb Z$-module inv. $T$, une $T$-structure à lacets
sur un groupe à lacets $\pi$ (avec $\pi\ne0$), muni
d'un isom $T\simeq T_\pi$, \deleted{et si $\Sigma$ opère sur
$T$ (\illegible{3} $\chi$), on dira qu'une action
est compatible}

On a des variantes profinies
(en particulier premières à $p$, où $p$ est premier
donné\dots) --- mais il y a dès mainte-
nant à signaler deux pts à vérifier dans
ce cas :
\[
\left\{
\begin{array}{l}
\operatorname{Norm}_{\pi}(L_i)=\operatorname{Centr}_{\pi}(L_i)=L_i
\quad\inserted{\text{dans le cas }\nu=1,\ \illegible{2}\text{ pas de pb}}\\[4pt]
\operatorname{Centr}(\pi)=(1)\quad(\text{cas anabélien})
\end{array}
\right.
\]
plus généralement, le centralisateur
de \unsure{tt} ss-groupe ouvert de $\pi$
est réduit à $1$\dots
\margin{\unsure{Ce n'est} \uline{démontré} pour le moment dans \uline{aucun} cas anabélien profini !}

\figure{[Trait isolé oblique descendant vers la droite, sous le texte.]}
\margin{117}

\scanpage{119}
\mspage{58}

14. Digression cohomologique (sur le ``bouchage des
trous).

Soit $X$ un schéma loc. noeth. \uline{régulier} de
\unsure{dim.} $1$, $S$ un sous-schéma fermé \inserted{réduit discret} tel que
$\forall s\in S$, $\dim_s S=1$ (donc $S$ défini par une
partie fermée discrète de $X$), $U=X\setminus S$.
\margin{Il vaudrait peut-être mieux démarrer avec les cas \unsure{purement topologiques}, \illegible{5} avec les groupes \uline{discrets}.}

On veut expliciter par voie galoisienne les
faisceaux \inserted{d'ens.} étales \inserted{constructibles} sur $X$ tels que $F\mid U$ soit
loc. ct. Par le topos d'Artin \inserted{\unsure{sur les ouverts de topos}}, ils corres-
pondent aux triples
\[
(F_U,F_S,\varphi)
\]
où $F_U$ est un faisceau \inserted{constr.} loc. constant sur $U$,
$F_S$ un faisceau (\unsure{néc.|nécess.} loc. constant)
et $\varphi$ un hom
\[
F_S\xrightarrow{\varphi}j^{*}i_{*}(F_U)
\]
(où $i:U\unsure{\hookrightarrow|\to}X$ et $j:S\to X$ sont les inclusions).
Toute condition de constructibilité étant mise à
part, un faisceau étale $F$ sur $X$ correspond
à un tel triple --- on se restreint ici aux
$F$ tels que $F_U$ \deleted{\unsure{provienne}} provienne du
topos fondamental $B_{\unsure{\Pi_1(U)|\pi_1(U)}}$ de $U$ \deleted{\illegible{1}}\dots
(Il n'y a pas lieu de supposer $F_U$ à fibres finies,

\margin{118}

\scanpage{120}

pour ce qui suit). \deleted{Nous allons}

Soit $\forall s\in S$ $\underline{O}_s$ un hensélisé de $\underline{O}_{X,s}$,
$D_s=\operatorname{Spec}\underline{O}_s$ \inserted{(``disque en $s$'')}, $D_s^{*}=D_s-\{s\}$ \inserted{$=\operatorname{Spec}K_s$ ($K_s$ corps des fractions de $\underline{O}_s$)} (``disque \unsure{pointé|ponctué}'' en
$s$) $D=\coprod_{s\in S}D_s$, $\unsure{D_s^{*}|D^{*}}=\coprod_{s\in S}D_s^{*}$. On a
un diagramme commutatif
\diagram{\begin{tikzcd}[ampersand replacement=\&]
\& D^{*} \arrow[dl, "\rho"'] \arrow[dr, "\sigma"] \& \& \\
U \arrow[dr, "i"'] \& \& D \arrow[dl, "\unsure{k|\kappa}"] \& S \arrow[l, dashed, "j_0"'] \arrow[dll, dashed, "j"] \\
\& X \& \&
\end{tikzcd}
[La branche pointillée $S\xrightarrow{j_0}D$ est placée entre parenthèses.]}
et on voit \inserted{(par Artin)} que ce diagramme \deleted{\illegible{2}} \inserted{permet} d'exprimer
\inserted{les faisceaux} \deleted{topos} étales de $X$ \deleted{\unsure{comme amalgamé}}
\deleted{des topos étales de $U,D$ sur $D^{*}$}
\deleted{\diagram{[Schéma biffé : objet supérieur \illegible{1}, flèches descendantes vers $U$ en bas à gauche et $D$ en bas à droite ; étiquettes \illegible{2}.]}}
comme des systèmes \deleted{de faisceaux} \inserted{$(F_U,F_D,\varphi)$}
avec $F_U$ faisceau étale sur $U$, $F_D$ faisceau
étale ``\unsure{ind} loc. constant'' sur $D$ [NB l'
inclusion $S\to D$ définit une équivalence
entre la cat. de ces faisceaux sur $D$, et
celle des faisceaux étales sur $S$] et
un hom. de faisceaux $F_D\to\sigma_{*}\rho^{*}F_U$ ou encore

\margin{119}
